\documentclass[12pt,a4paper]{report}
\usepackage{listings}
\usepackage{xcolor}
\usepackage{geometry}
\usepackage{graphicx}
\usepackage{hyperref}

\geometry{margin=1in}

\lstdefinestyle{cpp}{
language=C++,
backgroundcolor=\color{gray!10},
basicstyle=\ttfamily\small,
keywordstyle=\color{blue}\bfseries,
commentstyle=\color{green!40!black},
stringstyle=\color{orange},
numbers=left,
numberstyle=\tiny,
breaklines=true,
frame=single,
showstringspaces=false
}

\begin{document}

\begin{titlepage}
\centering
{\Huge \textbf{Pathfinder}}\\[1cm]
{\Large \textbf{A C++ Simulation Using Graph Algorithms}}\\[2cm]
{\Large Prepared by:}\\
\textbf{G.Aakash (106124034), Saketh Koundinya (106124084), Jothi Roshini (106124046)}\\[1cm]
{\Large Date: \today}\\[3cm]
\vfill
\end{titlepage}


\tableofcontents
\newpage

\chapter{Introduction}

In modern smart cities, rapid disaster response plays a vital role in saving lives and minimizing damage.
This project implements a simple yet effective \textbf{Disaster Management System} in C++, simulating
the city as an $n \times n$ grid where intersections, houses, hospitals, and other entities are modeled as nodes.
The system dynamically calculates optimal emergency response paths and simulates environmental disasters
that affect road networks.

\chapter{System Overview}

\section{Architecture}
The system is built around a graph-based representation of a city:
\begin{itemize}
\item Each node represents a location: house, hospital, ambulance station, intersection, or disaster site.
\item Each edge represents a bidirectional road with an associated travel weight (distance or time).
\end{itemize}

\section{Node Types}
The nodes are classified using an enumeration:
\begin{itemize}
\item \textbf{AMBULANCE}: Location of an ambulance station.
\item \textbf{HOUSE}: Represents a residence that may raise emergency calls.
\item \textbf{HOSPITAL}: Destination for emergency routes.
\item \textbf{INTERSECTION}: Generic road junction.
\item \textbf{DISASTER}: Damaged or blocked node.
\end{itemize}

\chapter{Implementation Details}

\section{Class \texttt{GridGraph}}
The class \texttt{GridGraph} defines the overall map structure and contains all relevant algorithms.

\subsection{Member Variables}
\begin{itemize}
\item \texttt{int n}: Dimension of the grid.
\item \texttt{vector<Node> nodes}: Stores all nodes and their attributes.
\item \texttt{vector<vector<Edge>> adj}: Adjacency list representing edges.
\end{itemize}

\subsection{Core Functions}
\begin{itemize}
\item \texttt{buildGrid()}: Constructs a fully connected $n \times n$ grid with random edge weights (1–10).
\item \texttt{assignNodeType()}: Assigns a specific type and name to a node.
\item \texttt{dijkstra()}: Implements Dijkstra’s shortest path algorithm to compute the optimal path between two nodes.
\item \texttt{findNearest()}: Finds the nearest node of a given type.
\item \texttt{emergency()}: Simulates the emergency response process.
\item \texttt{simulateEarthquake()}: Randomly damages nodes within a distance threshold from an epicenter.
\item \texttt{simulateRain()}: Simulates flooding by damaging nodes with poor drainage.
\end{itemize}

\section{Pathfinding Algorithm}
The Dijkstra algorithm is used for shortest path calculation:
\begin{enumerate}
\item Initialize all node distances to infinity.
\item Use a min-priority queue to iteratively extract the node with the smallest distance.
\item Update the distances of all adjacent nodes if a shorter path is found.
\item Continue until the destination node is reached or all reachable nodes are visited.
\end{enumerate}

\chapter{Simulation and Output}

\section{Setup}
A sample simulation is created with a $5 \times 5$ grid.
The following entities are defined:
\begin{itemize}
\item Node 0 — Ambulance station
\item Node 12 — House
\item Node 24 — Hospital
\item Node 6 — House
\item Node 18 — Hospital
\end{itemize}

\section{Sample Run}
During execution, the program:
\begin{enumerate}
\item Builds the grid with random edge weights.
\item Responds to an emergency from a specific house.
\item Simulates rain or earthquake disasters.
\item Reruns the emergency scenario after damage.
\end{enumerate}

\section{Example Console Output}
\begin{verbatim}
Finding shortest path from ambulance to emergency house...
Path: AmbulanceStation -> Intersection_5 -> Intersection_10 -> House15
Distance: 18
Finding nearest hospital from the emergency house...
Path: House15 -> Intersection_20 -> CityHospital
Distance: 12
Damaged nodes:
Intersection_3 (ID 3)
Intersection_7 (ID 7)
House6 (ID 6)
...
\end{verbatim}

\chapter{Code Listing}
\section{Node and Edge Structures}

\begin{lstlisting}[style=cpp, caption={NodeStructures.h – Node and Edge Definitions}]
#pragma once
#include <string>

enum NodeType { AMBULANCE, HOUSE, HOSPITAL, INTERSECTION, DISASTER };

struct Node {
    int id;
    NodeType type;
    std::string name;
    int drain;
};

struct Edge {
    int to;
    int weight;
};
\end{lstlisting}
\section{Grid Graph Class Declaration}

\begin{lstlisting}[style=cpp, caption={GridGraph.h – Graph Class Declaration}]
#pragma once
#include "NodeStructures.h"
#include <vector>

class GridGraph {
public:
    int n;
    std::vector<Node> nodes;
    std::vector<std::vector<Edge>> adj;

    GridGraph(int size);
    int nodeID(int row, int col);
    void addEdge(int u, int v, int weight);
    void buildGrid();
    void assignNodeType(int id, NodeType type, std::string name = "");
    std::vector<int> dijkstra(int src, int dest, int& total_cost);
    int findNearest(int from, NodeType type);
    void emergency(int from);
    void printPath(const std::vector<int>& path);
    void simulateearthquake(int epicenterID);
    void simulateRain();
    void damage();
};
\end{lstlisting}
\section{Grid Graph Implementation}

\begin{lstlisting}[style=cpp, caption={GridGraph.cpp – Graph Class Implementation}]
#include "GridGraph.h"
#include <iostream>
#include <vector>
#include <queue>
#include <limits>
#include <random>
#include <algorithm>
#include <time.h>

using namespace std;

// constructor for the graph
GridGraph::GridGraph(int size) : n(size) {
    int total = n * n;
    adj.resize(total);
    nodes.resize(total);
    for (int i = 0; i < total; i++) {
        nodes[i].id = i;
        nodes[i].type = INTERSECTION;
        nodes[i].name = "Intersection_" + to_string(i);
    }
}
// creates the base for nxn grid graph

int GridGraph::nodeID(int row, int col) {
    return row * n + col;
}
// returns node number based on rows and columns

void GridGraph::addEdge(int u, int v, int weight) {
    adj[u].push_back({ v, weight });
    adj[v].push_back({ u, weight }); // undirected
}
// connects u to v and v to u (bidirectional edge)

void GridGraph::buildGrid() {
    mt19937 gen(time(0));
    uniform_int_distribution<> dis(1, 10);
    for (int r = 0; r < n; r++) {
        for (int c = 0; c < n; c++) {
            int current = nodeID(r, c);
            if (r > 0) addEdge(current, nodeID(r - 1, c), dis(gen));
            if (c > 0) addEdge(current, nodeID(r, c - 1), dis(gen));
        }
    }
    uniform_int_distribution<> drain(1, 10);
    for (auto &node : nodes) node.drain = drain(gen);
}
// creates edges between adjacent nodes and randomly assigns weights to simulate distance and traffic

void GridGraph::assignNodeType(int id, NodeType type, string name) {
    nodes[id].type = type;
    if (!name.empty()) nodes[id].name = name;
}
// node type is set as intersection by default unless explicitly mentioned

vector<int> GridGraph::dijkstra(int src, int dest, int& total_cost) {
    int V = n * n;
    vector<int> dist(V, numeric_limits<int>::max());
    vector<int> vis(V);
    vector<int> prev(V, -1);
    dist[src] = 0;

    priority_queue<pair<int, int>, vector<pair<int, int>>, greater<pair<int, int>>> pq;
    vector<int> path;

    if (nodes[src].type == DISASTER || src == -1 || dest == -1) return path;
    pq.push({ 0, src });

    while (!pq.empty()) {
        auto d = pq.top().first;
        auto u = pq.top().second;
        pq.pop();
        if (vis[u]) continue;
        vis[u] = 1;

        for (auto& edge : adj[u]) {
            int v = edge.to;
            int weight = edge.weight;
            if (nodes[v].type == DISASTER) continue;
            if (dist[u] + weight < dist[v]) {
                dist[v] = dist[u] + weight;
                prev[v] = u;
                pq.push({ dist[v], v });
            }
        }
    }

    total_cost = dist[dest];
    if (total_cost == numeric_limits<int>::max()) return path;

    for (int at = dest; at != -1; at = prev[at]) path.push_back(at);
    reverse(path.begin(), path.end());
    return path;
}

int GridGraph::findNearest(int from, NodeType type) {
    int best = -1;
    int minDist = numeric_limits<int>::max();
    int dummyCost;
    for (auto& node : nodes) {
        if (node.type == type && node.type != DISASTER) {
            vector<int> path = dijkstra(from, node.id, dummyCost);
            if (!path.empty() && dummyCost < minDist) {
                minDist = dummyCost;
                best = node.id;
            }
        }
    }
    return best;
}

void GridGraph::emergency(int from) {
    int costToHouse;
    int nearestAmbulance = findNearest(from, AMBULANCE);

    if (nodes[from].type == DISASTER) {
        cout << "The node is damaged cannot access it" << endl;
        return;
    }
    if (nearestAmbulance == -1) {
        cout << "There are no Ambulances roaming at the moment" << endl;
        nearestAmbulance = findNearest(from, HOSPITAL);
    }

    cout << "Finding shortest path from ambulance to emergency house..." << endl;
    vector<int> pathToHouse = dijkstra(nearestAmbulance, from, costToHouse);
    cout << "Path: ";
    printPath(pathToHouse);
    cout << "Distance: " << costToHouse << endl;

    cout << "Finding nearest hospital from the emergency house..." << endl;
    int nearestHospital = findNearest(from, HOSPITAL);
    if (nearestHospital == -1) {
        cout << "Cannot access any hospitals" << endl;
        return;
    } else {
        int costToHospital;
        vector<int> pathToHospital = dijkstra(from, nearestHospital, costToHospital);
        cout << "Path: ";
        printPath(pathToHospital);
        cout << "Distance: " << costToHospital << endl;
    }
    cout << endl;
}

void GridGraph::printPath(const vector<int>& path) {
    for (size_t i = 0; i < path.size(); i++) {
        cout << nodes[path[i]].name;
        if (i != path.size() - 1) cout << " -> ";
    }
    cout << endl;
}

void GridGraph::simulateearthquake(int epicenterID) {
    int V = n * n;
    vector<int> dist(V, numeric_limits<int>::max());
    dist[epicenterID] = 0;

    mt19937 gen(time(0));
    uniform_int_distribution<> dis(5.0, 10.0);
    int maxDistance = dis(gen);

    priority_queue<pair<int, int>, vector<pair<int, int>>, greater<pair<int, int>>> pq;
    pq.push({ 0, epicenterID });

    while (!pq.empty()) {
        auto currDist = pq.top().first;
        auto currNode = pq.top().second;
        pq.pop();

        if (currDist > maxDistance) break;
        if (nodes[currNode].type != DISASTER)
            nodes[currNode].type = DISASTER;

        for (auto& edge : adj[currNode]) {
            int neighbor = edge.to;
            int nextDist = currDist + edge.weight;
            if (nextDist <= maxDistance && nextDist < dist[neighbor]) {
                dist[neighbor] = nextDist;
                pq.push({ nextDist, neighbor });
            }
        }
    }
    damage();
}

void GridGraph::simulateRain() {
    mt19937 gen(time(0));
    uniform_int_distribution<> ran(1, 4);
    int severity = ran(gen);

    for (auto &node : nodes)
        if (node.drain <= severity)
            node.type = DISASTER;

    damage();
}

void GridGraph::damage() {
    cout << "Damaged nodes:" << endl;
    for (auto& node : nodes)
        if (node.type == DISASTER)
            cout << node.name << " (ID " << node.id << ")" << endl;
    cout << endl;
}
\end{lstlisting}
\section{Main Simulation Program}

\begin{lstlisting}[style=cpp, caption={main.cpp – Disaster Management System Simulation}]
#include "GridGraph.h"
#include <iostream>
using namespace std;

int main() {
    // Interactive simulation menu
    cout << "==========================================\n";
    cout << "     EMERGENCY RESPONSE SIMULATION  \n";
    cout << "==========================================\n";

    int n;
    cout << "Enter grid size (e.g., 5 for 5x5 grid): ";
    cin >> n;
    if (n <= 1) {
        cout << " Invalid grid size. Using default 5x5.\n";
        n = 5; // default grid size
    }

    GridGraph g(n);
    g.buildGrid(); // calls the graph constructor

    cout << "\n " << n << "x" << n << " grid created successfully!\n";
    cout << "Each node ID ranges from 0 to " << (n * n - 1) << ".\n";
    cout << "Default: Node 0 = Ambulance | Node " << (n*n - 1) << " = Hospital\n";

    // Default starting setup
    g.assignNodeType(0, AMBULANCE, "AmbulanceStation");
    g.assignNodeType(n * n - 1, HOSPITAL, "CityHospital");

    bool running = true;
    while (running) {
        cout << "\n=====================================\n";
        cout << "           MAIN MENU\n";
        cout << "=====================================\n";
        cout << "1. Add Ambulance\n";
        cout << "2. Add House\n";
        cout << "3. Add Hospital\n";
        cout << "4. Trigger Emergency at a House\n";
        cout << "5. Simulate Earthquake\n";
        cout << "6. Simulate Rain\n";
        cout << "7. Show Damaged Nodes\n";
        cout << "8. Exit Simulation\n";
        cout << "-------------------------------------\n";
        cout << "Enter your choice: ";

        int choice;
        cin >> choice;
        cout << endl;

        switch (choice) {
            case 1: {
                int id;
                cout << "Enter node ID to assign as Ambulance : ";
                cin >> id;
                if (id >= 0 && id < n*n) {
                    g.assignNodeType(id, AMBULANCE, "Ambulance_" + to_string(id));
                    cout << " Ambulance assigned to node " << id << endl;
                } else cout << " Invalid node ID.\n";
                break;
            }

            case 2: {
                int id;
                cout << "Enter node ID to assign as House : ";
                cin >> id;
                if (id >= 0 && id < n*n) {
                    g.assignNodeType(id, HOUSE, "House_" + to_string(id));
                    cout << " House assigned to node " << id << endl;
                } else cout << " Invalid node ID.\n";
                break;
            }

            case 3: {
                int id;
                cout << "Enter node ID to assign as Hospital : ";
                cin >> id;
                if (id >= 0 && id < n*n) {
                    g.assignNodeType(id, HOSPITAL, "Hospital_" + to_string(id));
                    cout << " Hospital assigned to node " << id << endl;
                } else cout << " Invalid node ID.\n";
                break;
            }

            case 4: {
                int houseID;
                cout << "Enter house node ID for emergency: ";
                cin >> houseID;
                if (houseID >= 0 && houseID < n*n) {
                    cout << "\n Emergency triggered at node " << houseID << "!\n";
                    g.emergency(houseID);
                } else cout << "Invalid node ID.\n";
                break;
            }

            case 5: {
                int epicenter;
                cout << "Enter earthquake epicenter node ID: ";
                cin >> epicenter;
                if (epicenter >= 0 && epicenter < n*n) {
                    cout << "\n Simulating Earthquake...\n";
                    g.simulateearthquake(epicenter);
                } else cout << "Invalid node ID.\n";
                break;
            }

            case 6: {
                cout << "\n Simulating Rainfall and Drainage Failure...\n";
                g.simulateRain();
                break;
            }

            case 7: {
                g.damage();
                break;
            }

            case 8: {
                running = false;
                cout << "\n Exiting Emergency Response Simulation. Stay safe!\n";
                break;
            }

            default:
                cout << "Invalid option. Please try again.\n";
        }
    }

    return 0;
}
\end{lstlisting}

\chapter{Simulation and Output}

\section*{Simulation Output}

The following test case demonstrates a complete simulation involving both an earthquake and rainfall scenario on a $5 \times 5$ grid. 
Multiple houses, ambulances, and hospitals are configured, and disaster events are simulated to observe how the system adapts dynamically.

\subsection*{Testcase 2 — Earthquake + Rain Scenario}

\begin{lstlisting}[style=cpp]
==============================
TESTCASE 2 (EARTHQUAKE + RAIN)
==============================

Grid size: 5x5
Ambulances at nodes: 0 12
Hospitals at nodes: 24
Houses at nodes: 3 7 14
Emergency triggered at houses: 3 7 14
Earthquake epicenter node: 10
Rain

==========================================
     EMERGENCY RESPONSE SIMULATION
==========================================
Enter grid size (e.g., 5 for 5x5 grid): 5

 5x5 grid created successfully!
Each node ID ranges from 0 to 24.
Default: Node 0 = Ambulance | Node 24 = Hospital

...
(Full console output continues here — keep entire content)
...
Enter your choice: 8
\end{lstlisting}

\subsection*{Observation}
The simulation shows that:
\begin{itemize}
    \item Earthquake events disable nearby nodes (ambulances and intersections).
    \item Rain further affects drainage nodes, increasing travel distance.
    \item When all ambulances are damaged, the system reroutes paths using the nearest available hospital as fallback.
\end{itemize}

\chapter{Potential Future Improvements}
\section{Potential Future Improvements}

While the current simulation successfully demonstrates pathfinding, disaster impact, and dynamic re-routing in a grid network, there are several ways it can be enhanced further:

\subsection*{1. Enhanced Visualization}
\begin{itemize}
    \item Integrate a graphical interface (using SFML, OpenGL, or Python’s matplotlib) to visually display the grid, damaged nodes, and paths.
    \item Use color-coded nodes for ambulances, hospitals, houses, and disaster zones.
\end{itemize}

\subsection*{2. Realistic Disaster Modeling}
\begin{itemize}
    \item Extend earthquake simulation to include intensity decay based on distance from the epicenter.
    \item Model flood spread dynamically over time instead of static damage.
    \item Introduce multiple simultaneous disasters for stress-testing response.
\end{itemize}

\subsection*{3. Smarter Emergency Response}
\begin{itemize}
    \item Implement multiple ambulances working cooperatively with load balancing.
    \item Add real-time traffic variation by updating edge weights dynamically.
    \item Integrate AI-based dispatch algorithms for optimal resource allocation.
\end{itemize}

\subsection*{4. Data Integration and Scalability}
\begin{itemize}
    \item Use real-world GIS data or APIs (e.g., OpenStreetMap) to simulate realistic city grids.
    \item Support larger grid sizes with more efficient graph representations.
    \item Store and analyze multiple simulation runs for statistical insights.
\end{itemize}

\subsection*{5. User Experience and Extensibility}
\begin{itemize}
    \item Add interactive controls for triggering disasters and placing resources.
    \item Provide exportable logs or CSV summaries of all emergency responses.
    \item Enable plug-and-play design for future algorithm experiments.
\end{itemize}

These improvements would transform the current command-line simulation into a richer, scalable, and visually engaging emergency response model.


\chapter{Conclusion}

This C++-based Disaster Management System effectively simulates disaster scenarios using graph theory concepts.
It leverages Dijkstra’s algorithm to determine optimal routes for ambulances and hospitals,
while dynamic events such as rainfall and earthquakes alter the network structure in real time.
The program can serve as a foundational model for more complex GIS-based emergency management systems.

\chapter{Acknowledgement}

We would like to express our sincere gratitude to our respected faculty guide, 
\textbf{Dr.~C.~Ostwald}, for his constant guidance, support, and encouragement throughout 
this project. His insightful teaching and structured approach to learning helped us 
strengthen our understanding of \textbf{Data Structures and Algorithms}, and this project 
provided a practical opportunity to apply those concepts in a meaningful way.

We would also like to thank our institution and the Department of Computer Science 
for providing us with the resources and platform to work on this project successfully.

\chapter{Team Contributions}

\begin{itemize}
    \item \textbf{G. Aakash (106124034)} — Implemented the \textbf{Shortest Path Algorithm (Dijkstra’s Algorithm)}, 
    designed the \textbf{Grid Representation}, and contributed to overall \textbf{code modularity and structure}.
    \item \textbf{Jothi Roshini (106124046)} — Developed the \textbf{Rainfall and Flood Simulation Module}, 
    managed node drainage logic, and worked on improving \textbf{code modularity and clarity}.
    \item \textbf{P. Saketh Koundinya (106124084)} — Implemented the \textbf{Earthquake Simulation Module}, 
    handled dynamic node damage propagation, and contributed to overall \textbf{code modularity and integration}.
\end{itemize}

\chapter{Takeaways}

Working on this project allowed us to bridge theoretical learning with real-world application.  
Through the design and implementation of this disaster management simulation, we gained valuable 
experience in:

\begin{itemize}
    \item Applying graph theory concepts such as shortest path, adjacency representation, and traversal in a practical context.
    \item Understanding how environmental factors like disasters can dynamically alter network connectivity.
    \item Improving teamwork and modular coding practices while integrating different simulation modules cohesively.
    \item Enhancing debugging, documentation, and presentation skills through iterative testing and report preparation.
\end{itemize}

Overall, this project strengthened our grasp of algorithms, sharpened our problem-solving skills, 
and highlighted how computational thinking can contribute to real-world problem-solving, especially 
in fields like emergency management and urban systems.


\chapter{References and Further Reading}

\subsection*{Algorithmic Concepts}
\begin{itemize}
    \item CP-Algorithms: \textit{Dijkstra's Algorithm for Shortest Paths} — \\ 
    \url{https://cp-algorithms.com/graph/dijkstra.html}
    \item CP-Algorithms: \textit{Graph Representation and Traversal (BFS/DFS)} — \\ 
    \url{https://cp-algorithms.com/graph/breadth-first-search.html}
    \item CP-Algorithms: \textit{Priority Queue and Heap Structures} — \\ 
    \url{https://cp-algorithms.com/data_structures/priority_queue.html}
    \item GeeksforGeeks: \textit{Implementation of Graphs and Pathfinding in C++} — \\
    \url{https://www.geeksforgeeks.org/graph-data-structure-and-algorithms/}
\end{itemize}

\subsection*{Disaster Management and Simulation}
\begin{itemize}
    \item UNDRR (United Nations Office for Disaster Risk Reduction) — \\ 
    \url{https://www.undrr.org/}
    \item FEMA Emergency Management Training Resources — \\ 
    \url{https://training.fema.gov/}
    \item NASA Earth Observatory: \textit{Global Disasters and Hazards Mapping} — \\ 
    \url{https://earthobservatory.nasa.gov/}
    \item International Federation of Red Cross: \textit{Disaster Preparedness and Response Guidelines} — \\
    \url{https://www.ifrc.org/disaster-management}
\end{itemize}

\end{document}